   \documentclass[twocolumn,letterpaper]{article}

   % --- Essential Packages ---
   \usepackage[utf8]{inputenc}
   \usepackage[T1]{fontenc}
   \usepackage{amsmath, amssymb, amsfonts}
   \usepackage{geometry}
   \usepackage{graphicx}
   \usepackage{algorithm}
   \usepackage{algpseudocode}
   \usepackage{booktabs}
   \usepackage{xcolor}
   \usepackage{abstract}

   % --- Page Layout ---
   \geometry{left=15mm, right=15mm, top=20mm, bottom=20mm}

   % --- Title and Author Info ---
   \title{\textbf{Adaptive Stochastic Motion-Control and High-Frequency\\Safety Guardrails in Industrial Cutting Processes via Protocol-Agnostic Telemetry
 Ingestion}}
   \author{Automated Intelligence Research Division \\ \textit{Technical Specification Report}}
   \date{\today}

   \begin{document}

   \maketitle

   \begin{abstract}
   This paper presents a hierarchical control framework designed for high-speed industrial band saw operations, emphasizing machine adaptability and
 structural integrity. We propose a dual-layered intelligence architecture: a Stochastic Motion Control (SMC) subsystem that utilizes an ensemble of
 decision trees for real-time velocity optimization, and a High-Frequency Safety Guardrail providing deterministic response to mechanical impedance
 spikes. By utilizing a protocol-agnostic telemetry ingestion model, the system maintains high-fidelity operational awareness regardless of the underlying
 industrial communication standard. This document formalizes the mathematical landscapes governing the intelligence and provides the foundation for
 experimental validation.
   \end{abstract}

   % --- SECTION 1: INTRODUCTION ---
   \section{Introduction}
   In modern precision manufacturing, tool performance is highly sensitive to material heterogeneity. Traditional control-loop architectures rely on
 deterministic setpoints that fail to account for dynamic tool loading or stochastic variations in cutting resistance. This leads to suboptimal throughput
 and increased mechanical stress.

   To address these challenges, we introduce an intelligent motion framework. Unlike static controllers, our approach treats industrial movement as a
 dynamic optimization problem. By combining predictive ensemble learning with high-frequency reactive safety protocols, the machine can adapt its kinetic
 profile to real-time environmental variations while ensuring absolute mechanical protection through deterministic overrides.

   % --- SECTION 2: SYSTEM ARCHITECTURE ---
   \section{System Architecture}
   The architecture is organized into a tiered hierarchy designed to decouple intelligence from high-priority safety tasks.

\subsection{The Control Cycle}
The system operates on a strictly timed execution cycle ($\Delta t \approx 100$ms) for the primary control loop. However, the Safety Guardrail utilizes a high-priority hardware interrupt to ensure sub-millisecond response times, preempting all predictive optimization tasks in the event of a detected mechanical impulse. Each iteration consists of three phases:
   \begin{enumerate}
       \item \textbf{Observation:} Sensor telemetry is ingested and normalized via the industrial fieldbus interface.
       \item \textbf{Prediction:} The stochastic ensemble calculates optimal velocity offsets.
       \item \textbf{Execution/Intervention:} Commands are dispatched to actuators or high-priority safety overrides are triggered based on current load
 state.
   \end{enumerate}

   \subsection{Execution Priority Hierarchy}
   The system implements a "Safety-First" preemptive model. The Safety Guardrail operates at a higher priority level than the Stochastic Controller. In
 the event of a detected mechanical impulse, the safety layer preempts all predictive optimization tasks to ensure immediate response, minimizing tool
 damage and ensuring deterministic stability.

   % --- SECTION 3: DATA ACQUISITION & SENSOR MODEL ---
   \section{Hardware-to-Software Data Mapping}
   The fidelity of the control loop is dependent on the accuracy of the sensor state representation. The system ingests physical quantities through a
 specialized industrial interface.

   \subsection{Sensor State Vector}
   The machine's operational state is represented by a high-dimensional vector $\mathbf{S}$, comprising thermodynamic, mechanical, and kinematic
 observations. These include:
   \begin{itemize}
       \item \textbf{Kinematic Variables:} Band cutting velocity ($v_{\rightarrow}$), vertical descent speed ($v_{\downarrow}$), and lateral band
 deviation ($\delta_{band}$).
       \item \textbf{Mechanical Loading:} Motor current/amperage ($i$) and torque percentage ($\tau_{\%}$).
       \item \textbf{Environmental Context:} Ambient temperature ($T_{amb}$) and humidity ($\mathcal{H}_{rel}$).
   \end{itemize}

   \subsection{Protocol-Agnostic Telemetry Ingestion}
   To ensure cross-platform compatibility (e.g., Modbus TCP/IP, EtherNet/IP, Profinet), we utilize a high-level abstraction layer for data acquisition via
 an industrial fieldbus interface $\beta$. We implement a mapping function $\mathcal{M}$ that converts raw hardware-addressed identifiers $A_k$ into
 standardized engineering units $U_{eng}$:
   \begin{equation}
       u_i = \mathcal{M}(A_k) = \mathcal{F}(\text{raw}_k, \text{scale})
   \end{equation}

To mitigate high-frequency signal jitter inherent in industrial environments, all features $\xi$ undergo temporal smoothing via a sliding-window
  average. Note that increasing the window size $N$ improves signal fidelity but introduces significant phase lag (latency) into the control loop:
\begin{equation}
    \bar{\xi}_t = \frac{1}{N} \sum_{j=0}^{N-j} \xi_{t-j}
\end{equation}

   % --- SECTION 4: STOCHASTIC MOTION CONTROL ---
   \section{Stochastic Motion Control (SMC)}
   The SMC subsystem performs real-time velocity optimization based on the processed feature vector $X_t$.

   \subsection{Predictive Ensemble Optimization}
   We utilize a Bagging ensemble of $K$ decision trees. The system calculates a scaling coefficient $C_{final}$:
   \begin{equation}
       C_{final} = \text{clamp}\left(\frac{1}{K} \sum_{k=1}^{K} T_k(X_t), -1, 1\right) \cdot \lambda
   \end{equation}

   \subsection {Velocity Adaptation Law}
   The coefficient is applied to the vertical ($v_{\downarrow}$) and horizontal ($v_{\rightarrow}$) velocity components. To prevent actuator oscillation,
 a threshold $\epsilon$ is utilized:
   \begin{equation}
       \text{IF } |\sum_{i=1}^{n} \Delta v_i| > \epsilon \implies \text{Update Actuators}
   \end{equation}

   % --- SECTION 5: SAFETY GUARDRAIL ---
   \section{High-Frequency Safety Guardrail}
   The safety layer provides a deterministic response to mechanical impedance.

   \subsection{Impulse Detection Strategy}
   Using a history buffer $(H, \tau)$, the system calculates the interpolated baseline $\tau_{hist}$ at target depth $H_{target}$:
   \begin{equation}
       \tau(h) = \tau_1 + (h - h_1)\frac{\tau_2 - \tau_1}{h_2 - h_1}
   \end{equation}

\subsection {The Reactive Intervention Law}
A collision event is triggered if the impulse $\Phi$ exceeds the relative threshold $\Phi_{\text{threshold}}$ and a minimum absolute change $\epsilon_{min}$:
\begin{equation}
    \text{IF } \left(\frac{\tau_{current} - \tau_{hist}}{\tau_{hist}} > \Phi_{\text{threshold}}\right) \land (\tau_{current} - \tau_{hist} > \epsilon_{\text{min}}) \implies \text{Engage BRAKE\_STATE}
\end{equation}

   % --- SECTION 6: EXPERIMENTAL METHODOLOGY ---
   \section{Experimental Methodology}
   Testing will be conducted using a high-speed industrial band saw equipped with integrated torque sensors and encoders. We shall evaluate predictive
 accuracy across various density gradients (Softwood, Hardwood, and Metallic) to validate the adaptive intelligence of the SMC layer.

   % --- SECTION 7: CONCLUSION ---
   \section{Conclusion}
   This framework provides a robust mathematical foundation for intelligent motion control. By integrating probabilistic optimization with high-frequency
 safety oversight, we ensure both operational efficiency and mechanical integrity.

   \end{document}
